\section{Mục tiêu, lộ trình và bản đồ thực
hành}\label{mux1ee5c-tiuxeau-lux1ed9-truxecnh-vuxe0-bux1ea3n-ux111ux1ed3-thux1ef1c-huxe0nh}

Tutorial này phát triển đầy đủ 20 chủ đề trong slide \emph{Developing
Multiagent Systems: A Practical Tutorial with Python and PyTorch}, đồng
thời hoàn thành năm bài thực hành ở Unit 17. Bản mở rộng thêm hợp đồng
thông điệp, timeout thực sự, kiểm tra độ mới tồn kho, dữ liệu huấn luyện
có chia nhóm, kiểm thử tự động, ablation và bảng đo thực nghiệm.
\cite{slides}

Kết quả cuối là một trợ lý mua hàng cục bộ có CatalogAgent,
RecommendationAgent, InventoryAgent, PolicyAgent và Coordinator; một bộ
định tuyến intent bằng PyTorch đã train; một baseline một thành phần; và
bộ ca lỗi dùng để kiểm tra phối hợp. Dữ liệu catalog, kho và câu yêu cầu
đều \textbf{tổng hợp phục vụ học tập}. Không có thanh toán, đặt mua, dữ
liệu bệnh nhân hay kết nối dịch vụ ngoài.

Lộ trình đề xuất gồm ba buổi: buổi 1 học khái niệm và chạy luồng mua
hàng; buổi 2 đọc hợp đồng, thêm kiểm chứng và thử lỗi; buổi 3 train
classifier, chạy benchmark và phân tích kết quả. Người học nên chạy ví
dụ nhỏ trước rồi mới đọc các mở rộng. Các lệnh bên dưới chạy từ thư mục
\passthrough{\lstinline!solution/!}; script tự tìm thư mục kết quả theo
vị trí tệp nên cũng có thể gọi bằng đường dẫn tuyệt đối.

\subsection{Chuẩn bị môi trường và chạy
nhanh}\label{chuux1ea9n-bux1ecb-muxf4i-trux1b0ux1eddng-vuxe0-chux1ea1y-nhanh}

Mã đã được kiểm tra bằng Python 3.9.6, NumPy và PyTorch CPU trong máy
hiện tại; với môi trường mới nên dùng Python 3.10 trở lên có bản PyTorch
phù hợp. Hệ MAS cơ bản chỉ dùng thư viện chuẩn; classifier và đồ thị cần
các gói ghi trong \passthrough{\lstinline!requirements.txt!}. Các phiên
bản thực tế được lưu trong JSON và tệp môi trường, giúp người đọc phân
biệt môi trường đã thử với khuyến nghị cài đặt.

\begin{lstlisting}[language=sh]
python3 -m venv .venv
source .venv/bin/activate
python3 -m pip install -r requirements.txt
python3 tutorial/slide_baseline.py
python3 tutorial/mas_system.py
python3 tutorial/train_intent.py
python3 tutorial/evaluate_mas.py
python3 -m unittest discover -s tutorial -p 'test_*.py'
\end{lstlisting}

Các bước huấn luyện và đánh giá được tách tệp. Nếu chỉ muốn xem demo
MAS, chưa cần train classifier. Nếu muốn chạy B2 hoặc kiểm thử liên quan
router, cần tạo checkpoint trước, hoặc dùng checkpoint đã kèm trong
\passthrough{\lstinline!results/tutorial/!}. Khi dùng Jupyter, mở
notebook tutorial, chọn kernel cùng môi trường rồi Run All; đầu trang
notebook ghi rõ phạm vi thực thi.

\section{Unit 1 --- Từ mô hình AI đến hệ đa tác
tử}\label{unit-1-tux1eeb-muxf4-huxecnh-ai-ux111ux1ebfn-hux1ec7-ux111a-tuxe1c-tux1eed}

Một agent tiếp nhận thông tin, lựa chọn hành động và chịu trách nhiệm
một phần mục tiêu. MAS có nhiều thành viên với quyết định tương tác qua
môi trường hoặc thông điệp. Một nhóm hàm thuần túy không tự động chứng
minh cần MAS; phải có ranh giới nhiệm vụ có ý nghĩa, chính sách hoặc tài
nguyên riêng và cơ chế phối hợp có thể quan sát.

Trong ví dụ mua hàng, mô hình xếp hạng chỉ trả điểm. Hệ thông minh còn
cần truy hồi đúng sản phẩm, kiểm tra ngân sách, xác nhận kho và phản hồi
khi không đủ bằng chứng. Chia thành agent giúp xác định ai sở hữu mỗi
quyết định. Mã chạy chung một tiến trình; ``đa tác tử'' ở đây nói về
trách nhiệm và chính sách, không phải nhiều máy hay nhiều LLM.

Thực hành: chạy \passthrough{\lstinline!mas\_system.py!}, quan sát đầu
ra của kịch bản normal, timeout, malformed, stale và forged. Với normal,
danh sách là p2 giá 25 và p4 giá 55. Đây là kết quả của lọc dữ liệu, xếp
hạng và kiểm chứng kho, không có suy luận ngôn ngữ được giả lập. Ghi lại
thành phần nào tạo ra từng trường trong kết quả.

Câu hỏi tự kiểm tra: nếu thay tất cả agent bằng một hàm lớn nhưng giữ
hành vi, lợi ích và chi phí thiết kế thay đổi thế nào? Câu trả lời cần
nói đến khả năng thay thế thành phần, kiểm thử, logging và overhead,
thay vì khẳng định số lượng agent quyết định trí thông minh.

\section{Unit 2 --- Chọn bài toán và xác định lý do dùng
MAS}\label{unit-2-chux1ecdn-buxe0i-touxe1n-vuxe0-xuxe1c-ux111ux1ecbnh-luxfd-do-duxf9ng-mas}

Các miền e-commerce, tourism và healthcare support đều có thể phân rã
thành chuyên môn, nhưng yêu cầu dữ liệu và mức rủi ro khác nhau. Mua
hàng cần catalog, giá và tồn kho; lịch trình cần giờ mở cửa, thời gian
đi và chi phí; hỗ trợ hành chính y tế cần nội dung đã duyệt và lịch hẹn.
Cùng một sơ đồ agent không có nghĩa dùng chung tiêu chí nghiệm thu.

Tutorial chọn e-commerce vì có thể tạo một môi trường nhỏ với ground
truth rõ. Mỗi sản phẩm có ID, giá, category và tags; kho có số lượng
theo ID. Điều kiện thành công là trả đúng danh sách hợp lệ hoặc trả đúng
trạng thái không thể xác minh. Đây là lựa chọn giúp kiểm chứng đầu cuối
mà không cần API key.

Chỉ nên tách agent nếu ranh giới giúp quyết định, phân quyền hoặc thử
nghiệm. CatalogAgent và InventoryAgent dùng hai loại dữ kiện khác nhau;
PolicyAgent kiểm tra điều kiện bắt buộc; RecommendationAgent có thể thay
bằng learned ranker sau này. Với một catalog năm sản phẩm, MAS không có
lý do để được giả định nhanh hơn baseline đơn; benchmark sẽ đo thay vì
suy đoán.

\section{Unit 3 --- PEAS, tính hợp lý và đặc tả nhiệm
vụ}\label{unit-3-peas-tuxednh-hux1ee3p-luxfd-vuxe0-ux111ux1eb7c-tux1ea3-nhiux1ec7m-vux1ee5}

PEAS gồm Performance, Environment, Actuators và Sensors. Với trợ lý mua
hàng: Performance là đúng ràng buộc, truy hồi đúng và không khẳng định
tồn kho cũ; Environment là catalog và kho tổng hợp; Actuators là trả
danh sách hoặc trạng thái lỗi; Sensors là truy vấn, ngân sách và kết quả
agent. Không có actuator thanh toán.

Tính hợp lý là chọn hành động tốt theo thông tin và nguồn lực hiện có.
Khi kho timeout, ``không xác minh được'' là kết quả hợp lý dù không đưa
ra khuyến nghị hoàn chỉnh. Tạo một câu trả lời tự tin không làm nhiệm vụ
thành công nếu thiếu dữ kiện bắt buộc. Do đó success phải bao gồm cả
những ca từ chối hoặc fallback đúng.

{\def\LTcaptype{none} % do not increment counter
\begin{longtable}[]{@{}
  >{\raggedright\arraybackslash}p{(\linewidth - 8\tabcolsep) * \real{0.2000}}
  >{\raggedright\arraybackslash}p{(\linewidth - 8\tabcolsep) * \real{0.2000}}
  >{\raggedright\arraybackslash}p{(\linewidth - 8\tabcolsep) * \real{0.2000}}
  >{\raggedright\arraybackslash}p{(\linewidth - 8\tabcolsep) * \real{0.2000}}
  >{\raggedright\arraybackslash}p{(\linewidth - 8\tabcolsep) * \real{0.2000}}@{}}
\toprule\noalign{}
\begin{minipage}[b]{\linewidth}\raggedright
Thành viên
\end{minipage} & \begin{minipage}[b]{\linewidth}\raggedright
Đầu vào
\end{minipage} & \begin{minipage}[b]{\linewidth}\raggedright
Đầu ra
\end{minipage} & \begin{minipage}[b]{\linewidth}\raggedright
Quyền
\end{minipage} & \begin{minipage}[b]{\linewidth}\raggedright
Tiêu chí kiểm tra
\end{minipage} \\
\midrule\noalign{}
\endhead
\bottomrule\noalign{}
\endlastfoot
Catalog & Query & Danh sách ID & Đọc catalog & Tag khớp, query rỗng
không trả toàn bộ \\
Recommender & Candidate ID, budget & Tối đa 3 ID & Đọc thuộc tính & Chỉ
candidate hợp lệ, giá không vượt budget \\
Inventory & ID & Stock và tuổi dữ liệu & Đọc kho & Có đủ ID, stock không
âm, dữ liệu mới \\
Policy & ID, candidate, budget & ID đã xác nhận & Chỉ kiểm chứng & Không
cho ID lạ hoặc vượt giá \\
Coordinator & Yêu cầu & Status và items & Điều phối & Envelope đúng,
deadline, đầu ra được kiểm \\
\end{longtable}
}

Mỗi hàng là một hợp đồng thiết kế. Khi một agent xử lý thêm trách nhiệm
ngoài hàng của mình, nên giải thích vì sao và bổ sung test. Ví dụ
recommender có thể lọc budget để tiết kiệm, nhưng PolicyAgent vẫn kiểm
lại vì output thành phần có thể sai. Lặp kiểm tra ở ranh giới tin cậy là
chủ đích, không phải hai định nghĩa nghiệp vụ khác nhau.

\section{Unit 4 --- Môi trường, mục tiêu và kiến trúc
hybrid}\label{unit-4-muxf4i-trux1b0ux1eddng-mux1ee5c-tiuxeau-vuxe0-kiux1ebfn-truxfac-hybrid}

Môi trường catalog nhỏ gần như tĩnh trong một lần gọi, nhưng dữ liệu kho
có tuổi và có thể trở nên cũ. Một đề xuất dựa trên dữ liệu vừa đọc chỉ
có hiệu lực tại thời điểm đó. Dữ liệu cũ một giờ được tạo trong kịch bản
stale nhằm buộc hệ thống phân biệt ``có số lượng trong bộ nhớ'' với ``đã
xác minh hiện tại''.

Kiến trúc hybrid dùng luật cho ràng buộc cứng và mô hình học cho quyết
định dự đoán. Trong tutorial, intent classifier có thể chọn luồng, còn
budget và provenance không giao cho mạng tự quyết. Một lỗi phân loại cần
được ghi là lỗi định tuyến; một sản phẩm vượt budget cần bị chặn cho dù
score mô hình cao đến đâu.

Khi mở rộng, có thể dùng BDI để tổ chức niềm tin, mục tiêu và ý định,
hoặc utility để cân bằng sở thích mềm. Tuy nhiên, bất biến như ``không
trả ID không tồn tại'' vẫn cần kiểm chứng trực tiếp. Không nên biến điều
kiện bắt buộc thành một trọng số nhỏ trong score rồi hy vọng mô hình tự
ưu tiên.

\section{Unit 5 --- Kiến trúc giao tiếp và sơ đồ vận
hành}\label{unit-5-kiux1ebfn-truxfac-giao-tiux1ebfp-vuxe0-sux1a1-ux111ux1ed3-vux1eadn-huxe0nh}

\begin{center}
\begin{tikzpicture}[>=Stealth,node distance=8mm,
 b/.style={draw=accent,rounded corners,align=center,font=\small,minimum height=9mm}]
\node[b,minimum width=11cm] (c) {Coordinator: kiểm hợp đồng, deadline, trạng thái kết thúc};
\node[b,below=of c,xshift=-4.4cm] (a) {Catalog\\SEARCH};
\node[b,right=6mm of a] (r) {Recommender\\RANK};
\node[b,right=6mm of r] (p) {Policy\\VALIDATE};
\node[b,right=6mm of p] (i) {Inventory\\CHECK};
\draw[->] (c.south)-|(a.north);\draw[->] (c.south)--(r.north);
\draw[->] (c.south)--(p.north);\draw[->] (c.south)-|(i.north);
\draw[->] (a)--(r);\draw[->] (r)--(p);\draw[->] (p)--(i);
\end{tikzpicture}
\captionof{figure}{Kiến trúc cục bộ: các quyết định có vai trò riêng, kết quả đi qua coordinator.}
\end{center}

Thông điệp có task ID, sender, receiver, kind và payload. Task ID liên
kết các bước thuộc cùng yêu cầu; kind xác định loại thao tác; payload
chứa dữ liệu có cấu trúc. Timestamp cho độ mới được mô hình hóa bằng
tuổi dữ liệu kho, tránh lẫn với thời gian đo latency. Trace riêng lưu
thời lượng thực thi từng lần gọi.

\begin{lstlisting}[language=Python]
@dataclass(frozen=True)
class Message:
    task_id: str
    sender: str
    receiver: str
    kind: str
    payload: dict[str, Any]
\end{lstlisting}

Mỗi agent xác nhận người nhận và loại yêu cầu; coordinator kiểm sender,
receiver, task ID và loại phản hồi. Cùng một payload nhưng task ID sai
bị từ chối vì có thể là kết quả của công việc khác. Hệ hiện tại là
in-process nên chưa có chữ ký hay xác thực mạng; schema validation không
nên được mô tả thành cơ chế bảo mật phân tán hoàn chỉnh.

\section{Unit 6 --- Quy trình kỹ thuật từ yêu cầu đến kiểm
thử}\label{unit-6-quy-truxecnh-kux1ef9-thuux1eadt-tux1eeb-yuxeau-cux1ea7u-ux111ux1ebfn-kiux1ec3m-thux1eed}

Bắt đầu bằng câu người dùng: ``tìm đồ portable không quá 100''. Chuyển
thành điều kiện kiểm chứng: query không rỗng, chỉ trả sản phẩm từ
catalog, giá không quá 100, stock dương, dữ liệu kho không quá 60 giây.
Xác định trạng thái dừng cho không kết quả và lỗi. Sau đó mới phân vai
và viết message.

Bản end-to-end tối thiểu cần chạy được trước khi thêm ML. Nó cung cấp
nơi kiểm tra lỗi đầu vào, lọc dữ liệu và tính đúng của logic. Khi thêm
classifier, giữ nguyên bộ test để phát hiện mô hình mới làm hỏng ràng
buộc cũ. Tách các thay đổi giúp biết lỗi đến từ orchestration hay
learning.

Tệp \passthrough{\lstinline!mas\_traces.json!} là bằng chứng hữu ích khi
debug: mỗi case có kết quả cùng chuỗi lời gọi. Nếu kết quả trống, kiểm
tra đó là no\_results, timeout, availability\_unverified hay
contract\_error. Các trạng thái này có ý nghĩa khác nhau với người dùng
và không nên bị gom thành ``AI không trả lời''.

\section{Unit 7 --- Python cơ bản và hợp đồng đầu
vào}\label{unit-7-python-cux1a1-bux1ea3n-vuxe0-hux1ee3p-ux111ux1ed3ng-ux111ux1ea7u-vuxe0o}

Catalog, kho và sản phẩm được lưu trong bộ nhớ để ví dụ dễ chạy.
\passthrough{\lstinline!Product!} và \passthrough{\lstinline!Message!}
dùng dataclass đóng băng, giúp hạn chế sửa vô ý ở cấu trúc ngoài. Lưu ý
frozen dataclass không tự làm các dict lồng bên trong bất biến; code
tránh sửa payload sau khi gửi và không quảng cáo cấu trúc này là deep
immutable.

Coordinator kiểm query là chuỗi, budget là số hữu hạn và không âm. Các
ca NaN cần được chặn vì so sánh với NaN có hành vi dễ gây lọt ràng buộc.
Tệp test có cả budget âm và NaN. Một kiểm tra kiểu chung chung như ``có
thể chuyển thành float'' chưa đủ để xác nhận ý nghĩa nghiệp vụ.

Thực hành sửa lỗi: thay query bằng chuỗi rỗng và quan sát no\_results;
thay budget bằng -1 và NaN để quan sát invalid\_input. Sau đó đọc
\passthrough{\lstinline!Message.validate!} và
\passthrough{\lstinline!Coordinator.call!}. Chỉ mở rộng kiểu dữ liệu
được chấp nhận khi có lý do và test phù hợp.

\section{Unit 8 --- Catalog và
recommendation}\label{unit-8-catalog-vuxe0-recommendation}

CatalogAgent tìm chuỗi truy vấn trong tên, category hoặc tags của sản
phẩm. Truy vấn được strip và chuyển lower. Query rỗng không được xem là
khớp mọi chuỗi. Cách tìm này minh bạch, nhưng không hỗ trợ lỗi chính tả,
đồng nghĩa, truy vấn nhiều điều kiện hoặc ngữ nghĩa tiếng Việt; đó là
giới hạn có chủ đích của tầng retrieval.

RecommendationAgent lọc candidate có trong catalog và giá không quá
budget, sắp theo giá tăng rồi lấy tối đa ba sản phẩm. Nó không phải một
recommender cá nhân hóa được học từ người dùng. Tên ``recommendation''
mô tả chức năng đưa gợi ý, không hàm ý đã train mô hình xếp hạng.

Với query portable và budget 100, laptop 850 và tablet 150 bị loại;
mouse 25 và backpack 55 còn hợp lệ. Với query sports, giày có giá 70
nhưng stock 0 nên danh sách cuối trống. Hai ca này phân biệt lọc theo
giá và xác minh tồn kho, tránh nhầm retrieval có kết quả với nhiệm vụ đã
hoàn tất.

\section{Unit 9 --- Coordinator và bất biến đầu
ra}\label{unit-9-coordinator-vuxe0-bux1ea5t-biux1ebfn-ux111ux1ea7u-ra}

Coordinator thực hiện SEARCH → RANK → VALIDATE → CHECK. Kết quả được
kiểm ở mỗi ranh giới. Một message đúng kiểu nhưng sai task ID vẫn thất
bại. PolicyAgent đối chiếu cả catalog, candidate ban đầu và budget, giúp
chặn một ranker đưa vào ID tự tạo hoặc sản phẩm ngoài truy vấn.

Đầu ra cuối có \passthrough{\lstinline!status!} và
\passthrough{\lstinline!items!}; từng item có ID, giá lấy từ catalog và
availability. Chỉ trạng thái đã xác minh mới có item được gắn verified.
Nếu kho cũ hoặc lỗi, hệ không tái sử dụng một kết luận availability từ
bước trước. Đây là tiêu chí cụ thể có thể kiểm bằng failure injection.

Một vấn đề chưa được giải quyết trong prototype là race condition giữa
kiểm kho và mua thật. Ở đây không có thao tác mua nên kết quả chỉ là
snapshot gợi ý. Nếu thêm checkout, cần reservation, transaction,
idempotency và kiểm tra lại trước xác nhận. Không thể coi việc đã CHECK
là đã khóa được hàng.

\section{Unit 10 --- PyTorch: từ mạng chưa train đến classifier
thật}\label{unit-10-pytorch-tux1eeb-mux1ea1ng-chux1b0a-train-ux111ux1ebfn-classifier-thux1eadt}

Slide minh họa một mạng nhận hai đặc trưng số nhưng chưa huấn luyện. Bản
mở rộng thay bằng dữ liệu câu tiếng Việt có nhãn và bag-of-words. Quy
trình chuẩn hóa bỏ dấu, chuyển chữ thường và lấy token a--z; vocabulary
chỉ xây từ train. Vector đầu vào có 79 chiều trong lần chạy hiện tại,
mạng gồm Linear → ReLU → Linear với hidden size 32 và ba logits đầu ra.

Dữ liệu gồm 432 câu: ba lớp, mỗi lớp tám họ động từ, sáu cụm đối tượng
và ba hậu tố. Chia theo họ động từ: năm họ train, một validation và hai
test trong mỗi lớp; tương ứng 270/54/108 câu. Không có câu trùng chính
xác giữa các split và group ID không giao nhau. Tuy vậy, từ chỉ đối
tượng vẫn chia sẻ giữa split, nên đây là kiểm tra tổng hợp tương đối dễ.

Hàm mất mát CrossEntropyLoss nhận logits và nhãn số nguyên; không
softmax trước loss. Adam dùng learning rate 0.02, 120 epoch, seed 42 và
một CPU thread. Checkpoint được chọn theo validation loss thấp nhất;
test chỉ dùng để báo cáo cuối. Tài liệu PyTorch chính thức xác nhận giao
diện logits/target này. \cite{torchce}

\begin{lstlisting}[language=Python]
model.train()
optimizer.zero_grad()
logits = model(train_x)
loss = criterion(logits, train_y)
loss.backward()
optimizer.step()
model.eval()
with torch.inference_mode():
    probabilities = model(test_x).softmax(dim=1)
\end{lstlisting}

Classifier đạt accuracy và macro-F1 bằng 1 trên 108 câu test, mỗi lớp 36
câu. Đây là số đo thật trên tập tổng hợp, không chứng minh hiểu ngôn ngữ
rộng. Không có người dùng thật hoặc dữ liệu bệnh viện trong tập. Mục
tiêu của bài thực hành là thực hiện đúng vòng train--validation--test và
quan sát giới hạn ngoài phân phối.

\subsection{Ngưỡng tin cậy, từ chối và lỗi ngoài
miền}\label{ngux1b0ux1ee1ng-tin-cux1eady-tux1eeb-chux1ed1i-vuxe0-lux1ed7i-ngouxe0i-miux1ec1n}

Ngưỡng được chọn trên validation để giữ coverage cao với không có lỗi
quan sát được; kết quả là 0.5. Router còn yêu cầu câu có ít nhất một
token thuộc lexicon miền. Hai lớp kiểm tra này vẫn không bảo đảm phát
hiện OOD. Softmax là độ tự tin nội bộ của mô hình, không phải xác suất
được hiệu chỉnh rằng yêu cầu thuộc miền.

Sáu câu OOD minh họa được đánh giá sau chọn ngưỡng. Bốn câu được trả
clarify, nhưng ``thời tiết hôm nay thế nào'' bị route sang ecommerce và
``giải phương trình'' sang tourism. Nguyên nhân có thể liên quan token
trùng và độ tự tin cao ngoài miền; chưa có ablation đủ để kết luận định
lượng vai trò từng token. Báo cáo giữ nguyên các lỗi để cho thấy
accuracy test tổng hợp chưa đủ.

Một cải tiến tiếp theo là tạo tập development OOD và test OOD độc lập,
định nghĩa reject class hoặc huấn luyện phát hiện ngoài miền, rồi đo
false acceptance và coverage--risk. Không được sửa lexicon theo đúng sáu
câu đã xem rồi báo rằng hệ đã tổng quát hóa. Nếu dùng temperature
scaling, nhiệt độ cũng phải chọn bằng validation riêng.

\section{Unit 11 --- Timeout, fallback và giám
sát}\label{unit-11-timeout-fallback-vuxe0-giuxe1m-suxe1t}

\passthrough{\lstinline!asyncio.wait\_for!} áp deadline thực sự lên
coroutine gọi agent. Inventory chậm được mô phỏng bằng sleep 50 ms,
trong khi deadline mặc định là 10 ms. Khi quá hạn, lời gọi bị hủy và
coordinator trả timeout cùng danh sách trống. Đây là điểm mạnh hơn đoạn
\passthrough{\lstinline!try/except!} chỉ bắt TimeoutError mà không hề
thiết lập thời hạn.

Tính tương thích Python cần chú ý: ở một số phiên bản,
\passthrough{\lstinline!asyncio.TimeoutError!} và built-in
\passthrough{\lstinline!TimeoutError!} không hoàn toàn cùng lớp. Mã bắt
cả hai; test timeout đã được thực thi trong môi trường hiện tại. Trường
status của trace được khởi tạo trước try để ngay cả lỗi vẫn ghi nhận
được, tránh lỗi logging che mất lỗi gốc.

Deadline ở mức coroutine không tương đương khả năng dừng an toàn mọi thư
viện blocking hoặc một tiến trình ngoài. Với HTTP, cần timeout của
client; với CPU job, có thể cần process riêng; với thao tác ghi, timeout
còn để lại trạng thái không biết đã hoàn tất chưa. Prototype này chỉ đọc
dữ liệu nên không có retry giao dịch hoặc duplicate purchase.

Các chỉ số theo dõi gồm success, violation, P50/P95 latency và số
message. Một request--response được tính hai message. Trace hiện lưu
task ID và thời lượng, không chứa khóa API hay dữ liệu cá nhân. Cấu trúc
này đủ để minh họa quan sát vận hành mà chưa phải hệ tracing phân tán
đầy đủ.

\section{Unit 12 --- Case study thương mại điện
tử}\label{unit-12-case-study-thux1b0ux1a1ng-mux1ea1i-ux111iux1ec7n-tux1eed}

Thiết kế hoàn chỉnh có thể thêm intent, preference, catalog, ranker,
inventory và policy. Tuy nhiên, số lượng vai trò nên tăng khi có trách
nhiệm rõ. Với dữ liệu hiện tại, preference chưa có nguồn người dùng, vì
vậy không tạo hồ sơ giả rồi gọi là cá nhân hóa. Score giá tăng là một
baseline minh bạch.

Một bộ dữ liệu thật có thể đánh giá retrieval Recall@k, ranking NDCG@k
và tỷ lệ vi phạm điều kiện. Các chỉ số này cần relevance label hoặc
ground truth phù hợp. Bộ benchmark hiện dùng ID mong đợi và status,
không báo NDCG vì chưa có nhãn graded relevance. Phân biệt chỉ số chưa
đo giúp báo cáo giữ đúng phạm vi bằng chứng.

Gợi ý triển khai tiếp: thêm catalog adapter chỉ đọc, snapshot có version
và stock TTL theo loại hàng. Trước khi đánh giá, đóng băng tập yêu cầu
và catalog snapshot để mọi kiến trúc thấy cùng dữ liệu. Khi dùng thông
tin online, cần ghi timestamp để biết sai do mô hình hay do dữ liệu thay
đổi.

\section{Unit 13 --- Case study du
lịch}\label{unit-13-case-study-du-lux1ecbch}

Một tourism MAS có thể gồm preference, attraction, route và budget
agent. Trạng thái lịch trình chứa điểm đã chọn, thời điểm đến, thời gian
tham quan, chi phí và ràng buộc. RouteAgent phải tính cả di chuyển,
không chỉ cộng thời lượng tham quan. BudgetAgent phải bao gồm phí và vận
chuyển, không chỉ vé vào cửa.

Ví dụ tổng hợp: người dùng có 6 giờ và ngân sách 500 đơn vị; ba điểm cần
lần lượt 120, 90, 150 phút; tổng di chuyển 80 phút. Chọn cả ba cần 440
phút, vượt 360 phút dù từng điểm riêng lẻ đều phù hợp. Coordinator phải
bỏ điểm hoặc xin thay đổi ưu tiên mềm, không tạo lịch trình bất khả thi.

Bộ mã hiện không triển khai tourism planner và không dùng thông tin giờ
mở cửa thực tế. Đây là bài thiết kế chuyển miền hoàn chỉnh về vai trò và
ràng buộc, phù hợp với Unit 13; phần thực thi năm exercise được chọn ở
e-commerce. Nếu phát triển planner thật, dữ liệu mở cửa và thời gian di
chuyển phải lấy từ nguồn hiện hành có thể kiểm tra.

\section{Unit 14 --- Case study hỗ trợ hành chính y
tế}\label{unit-14-case-study-hux1ed7-trux1ee3-huxe0nh-chuxednh-y-tux1ebf}

Phạm vi phù hợp cho prototype sinh viên là chuẩn bị giấy tờ, hướng dẫn
quầy tiếp nhận và hỗ trợ lịch hẹn. IntakeAgent lấy yêu cầu;
InformationAgent chỉ đọc nguồn đã duyệt; SchedulingAgent đọc lịch;
PreparationAgent tạo danh sách việc cần chuẩn bị; Coordinator quản lý
quyền và xác nhận thao tác nếu có thay đổi lịch thật.

Bài này không triển khai phân loại cấp cứu, chẩn đoán hoặc điều trị. Lớp
hỗ trợ hành chính y tế trong classifier chỉ là nhãn tác vụ hành chính
tổng hợp. Do vậy accuracy của lớp này không được gọi là độ chính xác y
khoa. Từ ``healthcare'' trong một mô hình routing không tự mang ý nghĩa
clinical validation.

Đánh giá hành chính nên kiểm tính đúng thông tin, độ đầy đủ và bảo vệ dữ
liệu. Khi không có nguồn xác nhận, hệ phải trình bày chưa biết. Cách xử
lý inventory\_unverified ở e-commerce minh họa nguyên tắc tương tự,
nhưng không đủ để chứng nhận an toàn cho miền y tế.

\section{Unit 15 --- So sánh ba miền và quyền sở hữu ràng
buộc}\label{unit-15-so-suxe1nh-ba-miux1ec1n-vuxe0-quyux1ec1n-sux1edf-hux1eefu-ruxe0ng-buux1ed9c}

{\def\LTcaptype{none} % do not increment counter
\begin{longtable}[]{@{}
  >{\raggedright\arraybackslash}p{(\linewidth - 6\tabcolsep) * \real{0.2500}}
  >{\raggedright\arraybackslash}p{(\linewidth - 6\tabcolsep) * \real{0.2500}}
  >{\raggedright\arraybackslash}p{(\linewidth - 6\tabcolsep) * \real{0.2500}}
  >{\raggedright\arraybackslash}p{(\linewidth - 6\tabcolsep) * \real{0.2500}}@{}}
\toprule\noalign{}
\begin{minipage}[b]{\linewidth}\raggedright
Miền
\end{minipage} & \begin{minipage}[b]{\linewidth}\raggedright
Dữ liệu có thẩm quyền
\end{minipage} & \begin{minipage}[b]{\linewidth}\raggedright
Ràng buộc chính
\end{minipage} & \begin{minipage}[b]{\linewidth}\raggedright
Khi thiếu dữ liệu
\end{minipage} \\
\midrule\noalign{}
\endhead
\bottomrule\noalign{}
\endlastfoot
E-commerce & Catalog, price, inventory & Budget, ID hợp lệ, stock mới &
Không khẳng định availability \\
Tourism & Điểm đến, giờ mở cửa, thời gian đi & Khả thi thời gian, budget
& Trả phương án chưa xác minh hoặc sửa lịch \\
Healthcare administration & Tài liệu cơ sở và lịch hẹn & Đúng thông tin,
quyền, tối thiểu dữ liệu & Dừng thao tác và chuyển kiểm tra phù hợp \\
\end{longtable}
}

Ở cả ba miền, mỗi ràng buộc cần người sở hữu. Nếu ``ai cũng kiểm'' nhưng
không có hợp đồng, có thể không thành viên nào thực hiện đầy đủ.
Coordinator nên có một checklist kiểm được bằng máy cho điều kiện bắt
buộc, còn các đánh giá mềm để agent chuyên môn cung cấp.

Phân biệt hành động đề xuất với hành động ghi là cần thiết khi chuyển
prototype sang ứng dụng. Trong bài hiện tại, mọi actuator chỉ trả kết
quả nên có thể chạy tự động. Nếu thêm đặt hàng hoặc thay lịch, phải
thiết kế thao tác có khả năng xem trước và xác nhận rõ, cùng idempotency
và kiểm tra kết quả dịch vụ.

\section{Unit 16 --- Thiết kế benchmark có đối
chứng}\label{unit-16-thiux1ebft-kux1ebf-benchmark-cuxf3-ux111ux1ed1i-chux1ee9ng}

B0 là pipeline một thành phần có cùng kiểm chứng factual. B1 là MAS theo
luật có coordinator validation. B2 thêm classifier intent trước MAS. B3
tắt policy/freshness validation để đo ảnh hưởng của việc bỏ kiểm. Đây là
bốn cấu hình có mục đích khác nhau; B3 là ablation chẩn đoán, không phải
một sản phẩm được đề xuất sử dụng.

Tập đánh giá có 11 ca, mỗi ca lặp 20 lần, tổng 220 yêu cầu mỗi kiến
trúc. Các ca gồm bình thường, query không biết, query rỗng, budget 0,
hết hàng, budget âm, timeout, malformed, stale và forged. Repetition
dùng để ước lượng latency cục bộ; nó không biến 11 ca thành 220 tác vụ
độc lập về nội dung.

Success yêu cầu đồng thời status và danh sách ID bằng ground truth.
Violation kiểm item ngoài catalog, vượt budget hoặc khẳng định
availability từ dữ liệu stale. B0/B1/B2 đạt đủ các ca; B3 chỉ đạt 8/11
loại ca và có vi phạm. Số đo cụ thể được chèn từ CSV ở phần kết quả,
không viết tay.

Có hai giới hạn so sánh quan trọng. B0 mô phỏng trạng thái timeout đồng
bộ nên không chờ deadline thật; không thể dùng P95 của nó để kết luận
một hệ dịch vụ đơn luôn nhanh hơn MAS cùng tỷ lệ đó. B2 dùng cùng một
câu mua laptop để định tuyến trước mọi shopping case; benchmark này đo
overhead tích hợp classifier, \textbf{không} đo routing đa miền. Chất
lượng routing được đánh giá riêng trên tập 108 câu và OOD.

\section{Unit 17 --- Hoàn thành năm bài thực
hành}\label{unit-17-houxe0n-thuxe0nh-nux103m-buxe0i-thux1ef1c-huxe0nh}

\subsection{Exercise 1: Đặc tả tác tử và
PEAS}\label{exercise-1-ux111ux1eb7c-tux1ea3-tuxe1c-tux1eed-vuxe0-peas}

Ứng dụng chọn là shopping assistant. Bảng đặc tả ở Unit 3 có bốn worker
và coordinator, đầy đủ input, output, quyền và tiêu chí thành công. PEAS
bao gồm điều kiện từ chối đúng, không chỉ số lượng sản phẩm. Sơ đồ Unit
5 thể hiện nguồn quyết định và thứ tự truyền dữ liệu. Đó là bằng chứng
cho yêu cầu tối thiểu ba agent.

\subsection{Exercise 2: Mở rộng shopping
system}\label{exercise-2-mux1edf-rux1ed9ng-shopping-system}

Tệp \passthrough{\lstinline!slide\_baseline.py!} chép lại luồng bốn sản
phẩm từ Unit 7--9 và đã chạy năm kiểm tra: truy hồi tag, không vượt
budget, unknown query, empty query và chỉ trả candidate. Output được lưu
trong \passthrough{\lstinline!slide\_baseline.json!}. Từ baseline hai
agent đó, bản mở rộng bổ sung inventory và policy; B0 trong benchmark là
một baseline đơn có cùng điều kiện kiểm chứng để so công bằng, khác bản
slide ban đầu. Policy kiểm ID nằm trong catalog/candidate và không vượt
budget; inventory kiểm stock; coordinator kiểm response envelope và TTL.
Kịch bản portable/100 chứng minh không trả laptop vượt ngân sách, còn
sports/100 chứng minh không trả sản phẩm stock 0.

Để tự mở rộng, thêm PreferenceAgent nhận danh sách ID cùng sở thích có
schema. Hãy giữ budget ở lớp filter và dùng sở thích chỉ xếp hạng trong
miền hợp lệ. Viết test rằng một sản phẩm có score cao nhưng vượt budget
vẫn không xuất hiện. Đây là cách kiểm tra integration, thay vì chỉ thử
riêng hàm chấm điểm.

\subsection{Exercise 3: Huấn luyện intent
classifier}\label{exercise-3-huux1ea5n-luyux1ec7n-intent-classifier}

\passthrough{\lstinline!train\_intent.py!} tạo dữ liệu, chia nhóm, huấn
luyện, chọn checkpoint, chọn threshold và xuất per-class metrics.
\passthrough{\lstinline!intent\_model.pt!} chứa state\_dict, vocabulary
và threshold. Không tải một checkpoint không rõ nguồn. Sau chạy, xem
confusion matrix, loss train/validation và các câu OOD, rồi giải thích
vì sao kết quả 100\% vẫn không đủ cho triển khai.

Bài đã có accuracy, precision, recall và F1 từng lớp trên test, đúng yêu
cầu của slide. Phần confidence threshold là cơ chế abstention có thể
chạy, đồng thời báo cáo lỗi false acceptance. Muốn cải thiện nên tạo tập
dữ liệu phong phú hơn và giữ test độc lập sau khi chỉnh quy tắc.

\subsection{Exercise 4: Failure
injection}\label{exercise-4-failure-injection}

Bốn lỗi bắt buộc là timeout, empty result, malformed message và stale
availability đều có code và test. Bổ sung forged recommendation, sai
correlation ID, budget NaN và budget âm. Người học có thể chạy
\passthrough{\lstinline!test\_tutorial.py!} để xem ca nào thất bại khi
cố ý bỏ validation.

Ví dụ thực hành: đổi
\passthrough{\lstinline!Coordinator(validate=False)!} rồi chạy stale và
forged. Kết quả vẫn có thể mang status ok nhưng vi phạm factual. Sau đó
bật validation, quan sát availability\_unverified hoặc contract\_error.
Bài tập cho thấy đánh giá an toàn phải xét nội dung và provenance, không
chỉ việc hàm trả về mà không ném exception.

\subsection{Exercise 5: So sánh kiến
trúc}\label{exercise-5-so-suxe1nh-kiux1ebfn-truxfac}

Bảng benchmark đi kèm cho thấy B0 và B1 có cùng tỷ lệ success trên bộ
ca. Vì vậy kết quả \textbf{không chứng minh chia nhiều agent tăng chất
lượng} ở catalog toy. B1 có thêm message và latency; lợi ích là ranh
giới trách nhiệm, trace và khả năng thay thành phần. B3 cho thấy bỏ kiểm
chứng làm giảm chất lượng dù vẫn có nhiều agent.

Một so sánh tiếp theo nên tăng quy mô dữ liệu và độc lập nguồn lỗi, giữ
cùng tác vụ, lặp theo nhiều tải và dùng cùng cơ chế timeout cho cả
B0/B1. Nếu muốn chứng minh lợi ích chuyên môn hóa, cần tác vụ mà các
agent có dữ liệu hoặc công cụ khác nhau thực sự, và baseline đơn được
cấp nguồn lực tương xứng.

\section{Unit 18 --- Kiểm thử, rubric và lỗi thường
gặp}\label{unit-18-kiux1ec3m-thux1eed-rubric-vuxe0-lux1ed7i-thux1b0ux1eddng-gux1eb7p}

Có 17 test tự động cho tutorial, bao gồm đường bình thường, ngân sách,
empty/unknown, hết hàng, input sai, timeout, message malformed, stale,
forged, correlation, trace, chia nhóm dữ liệu, OOD đơn giản và tương
đương baseline. Test cuối đối chiếu 20 tổ hợp query/budget bình thường.
Các test lỗi kiểm hành vi tại ranh giới, không chỉ gọi lại công thức
trong implementation.

Rubric của slide được đối chiếu bằng sản phẩm: PEAS và vai trò có trong
báo cáo; nguồn Python và README có sẵn; timeout thực và validation có
test; classifier đã huấn luyện; benchmark có dữ liệu thô; phạm vi
synthetic được ghi rõ; demo và slide trình bày có kèm. Các yêu cầu này
được gắn vào manifest để người đọc kiểm tra nhanh.

Những lỗi cần tránh: gọi mạng chưa train là predictor hợp lệ; dùng cùng
tập để chọn ngưỡng và báo test; lấy output giả làm bằng chứng chạy; gọi
nhiều hàm là chuyên gia mà không có trách nhiệm; xem response ok là
thành công dù có ID bịa; và suy ra readiness triển khai từ vài dữ liệu
tổng hợp.

\section{Unit 19 --- Kiến trúc tham chiếu và hướng phát
triển}\label{unit-19-kiux1ebfn-truxfac-tham-chiux1ebfu-vuxe0-hux1b0ux1edbng-phuxe1t-triux1ec3n}

Nếu mở rộng thành dịch vụ, nên thêm adapter dữ liệu có version, cấu hình
timeout theo loại công cụ, retry cho thao tác đọc có giới hạn và logging
không chứa thông tin nhạy cảm. Những thao tác ghi cần idempotency key,
đọc lại xác nhận và mô hình trạng thái giao dịch. Không nên tự thêm một
vòng retry vô hạn khi không hiểu nguyên nhân lỗi.

Nếu mở rộng learned component, nên bắt đầu bằng ranking dataset có nhãn
và đánh giá riêng retrieval/ranking. Giữ invariant kiểm được bằng máy ở
coordinator. Nếu thêm LLM, coi output là proposal và giữ validator; lưu
model/prompt version; đo số token, độ trễ, groundedness và tỷ lệ lỗi
theo nhóm yêu cầu.

Đối với dữ liệu intent, bước cải tiến ưu tiên là dữ liệu tự nhiên, từ
vựng chưa thấy, câu đa ý định, phủ định, thiếu thông tin và ngoài miền.
So sánh classifier với keyword baseline trên cùng dữ liệu. Không nên
tăng mạng sâu trước khi kiểm tra bộ dữ liệu có thực sự yêu cầu năng lực
đó hay chưa.

\section{Unit 20 --- Tổng kết, demo và câu hỏi bảo
vệ}\label{unit-20-tux1ed5ng-kux1ebft-demo-vuxe0-cuxe2u-hux1ecfi-bux1ea3o-vux1ec7}

Một bản demo 5--7 phút có thể đi qua sáu bước: giới thiệu PEAS; chạy
portable/100; mở trace; thử timeout và stale; trình bày confusion matrix
cùng lỗi OOD; kết thúc bằng bảng B0--B3. Đây là trình tự cho thấy cả cơ
chế và giới hạn, không chỉ một output thành công được chọn sẵn.

Câu hỏi bảo vệ gợi ý: vì sao dùng MAS ở catalog nhỏ; timeout nào thực sự
được áp dụng; vì sao softmax cao vẫn sai; classifier có phải RL không;
B3 đo điều gì; vì sao baseline đơn vẫn cần kiểm kho; và dữ liệu nào sẽ
cần thêm trước khi dùng thật. Các câu trả lời phải gắn với code, dữ liệu
và protocol đã công bố.

Thông điệp kỹ thuật của tutorial là xây một hệ có hợp đồng, kiểm chứng
và bằng chứng đo được. Phần học máy đóng góp một thành phần dự đoán đã
train và có giới hạn được đo; phần điều phối giữ các điều kiện quan
trọng bất kể dự đoán. Cả hai cần được kiểm tra bằng cùng tiêu chí nhiệm
vụ.
